% ---------------------------------------------------------------------
%  Tartalmi összefoglaló (magyar) és Abstract (angol)
% ---------------------------------------------------------------------
\chapter*{Tartalmi összefoglaló}
\addcontentsline{toc}{chapter}{Tartalmi összefoglaló}
\thispagestyle{plain}

A feleletválasztós vizsgakérdések írása lassú és hibalehetőségekkel teli szakmai munka: a kérdésnek
egyértelműnek kell lennie, pontosan egy válaszlehetőségnek kell a tananyag szerint helyesnek lennie, a
disztraktoroknak pedig hihetőnek, de igazolhatóan hibásnak. A nagy nyelvi modellek (LLM) megjelenésével
a legelterjedtebb gyakorlat az lett, hogy az oktató feltölti a jegyzetét egy felhőben futó csevegő
asszisztensbe, és „tíz feleletválasztós kérdést” kér. Ez kényelmes, de négy lényeges hiányossága van:
(P1) semmi sem ellenőrzi, hogy a megjelölt helyes válasz valóban következik-e a dokumentumból;
(P2) a kérdések száma, típusa és nehézsége csak kérés, a hosszú bemenetet a modellek egyenetlenül
használják, így a kérdések néhány szakaszra tömörülnek; (P3) a tananyag egy harmadik félhez kerül;
(P4) a módszer nagy, távoli modellekre és hosszú kontextusra támaszkodik.

Dolgozatunkban bemutatjuk a \Mimir{} rendszert, amely a fenti egyetlen promptot explicit, ellenőrizhető
\emph{módszerrel} váltja fel, és egyetlen fogyasztói GPU-n futó, 7 milliárd paraméteres, 4 bitre kvantált
nyelvi modellel működik. A módszer lépései: szemantikus darabolás ablakozott mondatbeágyazással;
lefedettségvezérelt \emph{vizsgaterv} (blueprint), amely a fogalmakat, a kérdéstípusokat és a Bloom-szinteket
a legnagyobb maradék elvén osztja ki a slotokra; slotonkénti visszakeresés (sűrű vagy BM25\,+\,sűrű,
reciprok rangfúzióval); egyenként generált, forrásrészletekre hivatkozó kérdések; valamint a legolcsóbbtól
a legdrágábbig rendezett ellenőrző lánc (forma, meta-hivatkozás, hivatkozás, alátámasztottság és
disztraktorok, vak megoldhatóság, válaszszivárgás, duplikátum), amely hiba esetén visszacsatolással
újrapróbálkozik. Egy fogalomgráf-alapú változat (GraphRAG-lite) a tervezőt, a disztraktorjavaslatokat
és a kontextust gazdagítja. A rendszer mikroszolgáltatás-architektúrája zéró adatmegőrzést valósít meg:
a feltöltött tartalom csak a feladat idejére létezik, a naplóban csak kriptográfiai ujjlenyomatok maradnak.

Második hozzájárulásunk egy reprodukálható, nyílt kiértékelési \emph{módszertan}: tizenöt, egyszerre
egy tényezőt változtató kísérleti kar; 39 magyar és angol dokumentumból, 326 referenciakérdésből álló
adatkészlet; a generátoroktól eltérő modellcsaládba tartozó LLM-bíró; formálisan definiált mutatók
(formai megfelelés, alátámasztottság, vak megoldhatóság, disztraktor-validitás, lefedettség, szivárgás);
valamint dokumentumszintű statisztika (bootstrap konfidenciaintervallum, páros Wilcoxon-próba
Holm-korrekcióval, rangbiszeriális hatásméret, Krippendorff-$\alpha$).

Egy magyar dokumentumon végzett pilot (tíz kar, három ismétlés, 300 kérdés) szerint a helyi modell a
GPU-memóriából 5,4--5,9\,GB-ot foglal, tehát 8\,GB-os kártyán is fut. A pilot két, korábbi méréseket
érvénytelenítő mérési hibát és egy új hibamódot is feltárt: az ellenőrző visszacsatolási hurka a helyes
választ a kérdés szövegébe szivárogtatta, amit a bíró sem büntetett. A javított csővezetékkel végzett fő
vizsgálat (14 kar, kilenc magyar és angol dokumentum, karonként hat, 1\,670 kérdés) szerint a helyi 7B
modellel futó, tervezett és ellenőrzött csővezeték érte el a legmagasabb minőségi indexet (0,96; 95\,\%-os
CI: 0,93--0,98), szemben a naiv RAG (0,64), a teljes dokumentumot egy promptban kapó helyi modell (0,83) és
a teljes dokumentumot kapó nagy szervermodell (0,95) értékével. A tervezés 0,24-dal, az ellenőrzés -- a
pilottal ellentétben -- mind a hat dokumentumon 0,08-dal javított; a más modellcsaládba tartozó ellenőrző és
a fogalomgráf nem javított; a szivárgási szabályok a szivárgást 0--3\,\%-ra csökkentették; ugyanaz a
csővezeték a szervermodellel 0,99-et ért el. Az ellenőrző ítélete a bíróéval csak gyengén egyezik
(Cohen-$\kappa\le0{,}22$), a naiv RAG minőségét pedig a darabolás apró részletei is meghatározták. Hat
dokumentumon a különbségek a Holm-korrekció után nem szignifikánsak; a teljes, 39 dokumentumos futás és a
bíró oktatói validálása a következő lépés.

\medskip
\noindent\textbf{Kulcsszavak:} automatikus kérdésgenerálás, visszakereséssel kiegészített generálás (RAG),
kis nyelvi modellek, verifikáció, LLM-bíró, adatvédelem, oktatástechnológia, magyar nyelv

% ---------------------------------------------------------------------
\chapter*{Abstract}
\addcontentsline{toc}{chapter}{Abstract}
\thispagestyle{plain}

\begin{otherlanguage}{english}
The most common way to obtain an exam from teaching material today is to upload the document to a
hosted chat-based large language model (LLM) and ask for a test. This offers no evidence that the keyed
answers are supported by the document, no control over which parts of the material are covered, and it
discloses the material to a third party. We present \Mimir{}, an open system that generates
multiple-choice exams from a teacher's own documents with a 4-bit, 7-billion-parameter model running on a
single consumer GPU. \Mimir{} replaces the single prompt by an explicit method: semantic chunking with
windowed sentence encoding, a coverage-driven exam \emph{blueprint} that allocates concepts, question
types and Bloom levels to slots by the largest-remainder rule, per-slot dense or hybrid (BM25 + dense,
reciprocal rank fusion) retrieval, per-question generation with chunk citations, and a cheapest-first
verifier chain with feedback-driven retries, optionally enriched by a session-scoped concept graph. The
microservice architecture retains no uploaded content beyond the lifetime of a job.

We also contribute an evaluation methodology and open harness: fifteen single-factor experiment arms, a
39-document Hungarian--English dataset with 326 reference questions, a cross-family LLM judge, formally
defined measures and document-level statistics. A pilot on one Hungarian document (ten arms, three seeds,
300 questions) showed that the local model occupies 5.4--5.9\,GB of GPU memory, exposed two measurement
defects, and revealed an answer-leakage failure mode that inflates both automatic and LLM-judge scores. In
the main study with the corrected pipeline (14 arms, nine Hungarian and English documents, six per arm,
1,670 questions), the planned and verified pipeline on the local 7B model reached the highest quality index
(0.96, 95\,\% CI 0.93--0.98), against 0.64 for naive RAG, 0.83 for the whole document in one prompt to the
same model and 0.95 for the whole document to a large server model. Planning added 0.24; verification,
unlike in the pilot, added 0.08 on all six documents; a verifier from another model family and the concept
graph did not help; leakage rules cut leakage to 0--3\,\%; and the same pipeline on the server model reached
0.99. The verifier agrees only weakly with the judge (Cohen's $\kappa\le0.22$), and naive RAG proved fragile
to chunking details. With six documents no difference survives Holm correction; the full 39-document run and
the teachers' validation of the judge are the next step.

\medskip
\noindent\textbf{Keywords:} automatic question generation, retrieval-augmented generation, small language
models, verification, LLM-as-a-judge, privacy, educational technology, Hungarian
\end{otherlanguage}
